\section{Catalogue of excluded schedule families}\label{app:catalogue}

This appendix lists the explicit oblivious linear-time schedule families excluded by the results of this paper. A schedule family
prescribes, for each size $d$ in its domain, the motion of the input head and of the work heads and the append steps; a family is
\emph{excluded} if for every $C\ge1$ and every value width $w$ there is $d_0$ such that, at every size $d\ge d_0$ in its domain, no machine of
value width at most $w$ whose trace is the prescribed one transposes within $C(L+1)$ steps.

\paragraph{Status.} For rows 1--20 the excluding theorem is formal, and that the schedule belongs to the class is checked by inspection
of the schedule. \Cref{cor:H} and \cref{thm:Q} exclude size by size, with thresholds that depend only on the bounds for the number of tapes and the
value width (uniform over all machines within those bounds),
so the rows they cover are excluded in the sense above. \Cref{thm:BM,thm:J} speak of one machine at all large sizes: for the rows that
only they cover (rows 3--7 and 12) they exclude the machines that follow the schedule at all large sizes of its domain, not, size by size,
every machine that follows it at one size. For the schedules defined only for even $d$ or for powers of two (rows 3, 5 and 6) they are
applied along those sizes, which their formal class predicates, quantified over all large $d$, do not literally cover. For rows 21--38 the certificate is
built from formal lemmas of \cref{sec:certificates}, but its verification on the schedule is a written argument; each is stated below as a
proof sketch only. \emph{None of these exclusions is a numbered result of this paper.}

\paragraph{Crossings (numerical).} The last column records whether \cref{thm:X} alone could exclude the family. ``Yes'' means that the
schedule's work heads stay within $O(b)$ cells, so it has no crossing revisits of linear excursion at all. The other entries are
numerical, from exact computations at $d\le128$: ``no'' when the minimum set separating a large output set from the crossing steps of
revisits of excursion more than $2W$ had at least $0.25$ times as many elements as the output set; ``small'' when that ratio was at most
$0.07$, which suggests, as a heuristic, that \cref{thm:X} applies; and ``---'' when nothing was measured.

{\small
\begin{longtable}{@{}r p{0.42\textwidth} c p{0.2\textwidth} c@{}}
\toprule
& Schedule & $k$ & Excluded by & Crossings \\
\midrule
\endhead
1 & Direct copy: output bit $\ell$ emitted while reading input bit $\ell$; work tape idle (computes the identity) & 1 & \cref{cor:H}, \cref{thm:BM} & yes \\
2 & Persistent triangles: output $(j,i,s)$ emitted while reading input $(j,-j-i\bmod d,s)$; work head within $b$ cells & 1 & \cref{cor:H}, \cref{thm:J} & yes \\
3 & One-pass shear with row rotations, two banks and a row buffer; seven monotone phases ($d$ even) & 3 & \cref{thm:BM,thm:J} & no \\
4 & Ruler schedule: two banks moving one row per episode, a third head making excursions of $2$-adic lengths inside one row & 3 & \cref{thm:J} & --- \\
5 & Row expander: rows mixed in a two-head deep-nesting workspace of width $W$, two banks, four monotone phases ($d$ even) & 4 & \cref{thm:BM} & --- \\
6 & Fan-out cascade fed by fresh input slices at each scale ($d=2^m$) & 2 & \cref{thm:J} (suffix form) & --- \\
7 & Unit-slope passes: load, reflection, shifted sweep, emit; at most three constant-velocity segments per pass & 2 & \cref{thm:J} & small \\
8 & Bounce: zig-zag input, one tape sweeping $[0,L)$, nested excursions of height at most $L/(4d)$ on the other & 2 & \cref{thm:J}, \cref{cor:H} & small \\
9 & Sublinear stack: blocks of $L/(4d)$ outputs pushed and popped on fresh cells while the other tape sweeps & 2 & \cref{thm:J}, \cref{cor:H} & --- \\
10--11 & Free pseudo-random Dyck walk of length $L$ on one tape while the other sweeps the tape six times (two variants) & 2 & \cref{thm:J}, \cref{cor:H} & --- \\
12 & One push and one pop of height $2L$ on a stack tape while the other tape bounces & 2 & \cref{thm:J} & --- \\
13--16 & One free head with all other excursions $O(L/\log d)$: slot cascades and skewed wavefronts (four variants) & 2 & \cref{cor:H} & small (16) \\
17--20 & No free head, owner load at most $4$, family heights in at most $4$ scale windows: multi-scale push--pop, two-excursion hub, LIFO Dyck stack, torus shifts & 2 & \cref{thm:Q} & --- \\
21--23 & One push--pop excursion on a stack tape, with a free tape running a multi-scale cascade or bouncing in a hub (three variants) & 2 & slab-chunk genie on a two-scale grid & --- \\
24 & Multi-scale Dyck stack: $d$ sawtooth waves on a stack tape while the free tape runs a cascade and bounces in a hub ($d=2^\kappa$) & 2 & block cut on the hub & no \\
25--26 & Two push--pop excursions on a stack tape, with a mandatory cut at the second excursion (two variants) & 2 & block cut & --- \\
27 & Two-exception inverse-square geometry, swept variant~\cite{GapAreaMemo2026} ($d\ge64$) & 10 & closed post-region, $O(d^{1/3})$ boundary steps & no \\
28--29 & Two-exception inverse-square geometries~\cite{GapAreaMemo2026} ($d\ge64$) & 9, 11 & window lemma after two cut steps & small \\
30 & Mirrored dyadic fold cascade & 2 & window lemma (\cref{lem:stepinv}), uncut & small \\
31 & Unit-slope dyadic cascade ($d=2^m$) & 2 & sub-row window, uncut & --- \\
32 & Dyadic fold cascade with row blocks of doubling width & 2 & window lemma, uncut & no \\
33--34 & Two stack excursions with the free tape in a hub of width $L/2^s$ ($s$ fixed) or $d$ & 2 & block cut; space--time potential & --- \\
35 & Daisy-chained stack copies with golden-ratio shifts & 2 & piecewise-isometric windows & --- \\
36 & LIFO fan-out cascades of $\log d$ levels (two variants; $d=2^\kappa$) & 2 & laminar cut of $2^{1-p}L+3$ steps, $p=O(\log\wval)$ & no \\
37 & Equal-volume tapped cascade with $\log d$ meso-scale stages & 2 & meso-hub cut of $O(L/\log d)$ steps & no \\
38 & Two-band meso cascade & 2 & closed post-region, $O(L/\log d)$ boundary steps & --- \\
\bottomrule
\end{longtable}
}

\paragraph{Rows 1--20.} For bounded macro-phases (\cref{thm:BM,thm:J}) the crux rate is $1/(4(\abs B+1)^{M_{\mathrm{ph}}-1})$ with the bank $B$ and
the number of phases read off the schedule: for example, row~3 has banks $\{0,1\}$ and seven phases, row~5 banks $\{0,1\}$ and four phases.
Rows 13--16 satisfy the hypotheses of \cref{cor:H} with the free head named and all other excursions $o(L)$; rows 17--20 those of
\cref{thm:Q} with owner load and window counts read off the schedule. For row~1, \cref{thm:BM} gives crux rate $\frac14$; the true rate is at
most $\frac12$, because pairing $q$ with $\tpos(q)$ gives an exclusive-or code of $L/2$ bits.

\paragraph{Rows 21--38.} Proof sketch only. Rows 21--23 are excluded by the frozen-trace count used in the proof of \cref{thm:H}: on the pop
leg of the stack, revisits that span the peak are order-reversing, and a two-scale slab grid with a greedy slicing gives a trapped chunk whose
genie costs less than its output surplus. Rows 24--38 are excluded through \cref{thm:crux-excludes} by an explicit certificate. Rows 24--26 and 33:
the cut at a mandatory part of the stack motion makes every relevant output carry at most $a$ blocks of potential sources and every input
at most $e$ blocks, and \cref{lem:blockcut} gives an acyclic set of density $1/\binom{a+e}{a}$, uniformly in the depth. Rows 27--30 and 32: the
window lemma (\cref{lem:window,lem:stepinv}) with potentials $\varphi(q)=\tpos(q)$ and $\psi(p)=L-1-p$ or a row-local variant, possibly after
cutting a sublinear boundary. The geometries of rows 27--29 have hop counts of order $L/i^2$ in their $i$-th family; the gap-area memo of the
ninth paper shows that its frozen-gap-area certificate cannot exclude this case~\cite{GapAreaMemo2026}, whereas the certificates here use cut
steps and a window potential instead. Row 34 uses a space--time potential, row 35 a bound for piecewise-isometric motion with boundedly many
segments, row 36 a bounded-level laminar cut, row 37 a meso-scale hub cut, and row 38 the formal closed-post-region criterion
(\lean{isAcyclicIn_ballAvoidSystem_of_closedPostRegion}).

As a numerical check on the method, the same classification was run on explicit transposers that run in time $\omega(L)$ (shuffle and
butterfly transposers, a radix transposer on a stack, and one-tape transposers); none of the exclusions was triggered by them.
